En este informe se presentan la implementación y los resultados del análisis experimental de distintos algoritmos: Merge Sort, Quick Sort, Patience Sort y \texttt{std::sort} para el problema de ordenamiento de arreglos, y Naive y Strassen para la multiplicación de matrices cuadradas. El objetivo es contrastar las complejidades teóricas ya conocidas con el comportamiento real de cada algoritmo, midiendo tiempo de ejecución y uso de memoria bajo distintas condiciones de entrada. Para ello, los experimentos se ejecutaron en un entorno controlado (Intel Core i5-14400F, 16 GB RAM, C++17 compilado con GCC 16.2.1 y \texttt{-O3}), evaluando tamaños de hasta $10^7$ elementos y matrices de hasta $1024 \times 1024$, con distintos estados iniciales, tipos de matriz y dominios numéricos. Los resultados muestran que Naive es sistemáticamente más rápido que Strassen hasta $N=1024$ pese a su menor complejidad teórica, que \texttt{std::sort} y Merge Sort mantienen una estricta estabilidad en $O(n \log n)$, y que Quick Sort y Patience Sort se degradan drásticamente ante dominios numéricos con alta tasa de duplicados y, en el caso de Patience Sort, ante distribuciones ascendentes. Se concluye que las cotas teóricas deben contrastarse con los factores constantes de implementación en escenarios reales, donde la sobrecarga recursiva, el consumo de memoria y la estructura de los datos de entrada resultan determinantes para el rendimiento observado.
